\documentclass{article}
%%%%% PROGRAMMING PACKAGES
\usepackage{expl3,xparse}     %LaTeX3
\usepackage{mfirstuc}    %first letter uppercase with \makefirstuc (used in todo notes)
\usepackage{settobox}     %get box dimensions with \settoboxwidth (used in fit environment) => pretty stupid package.. could be avoided

%%%%% FIGURES, BOXES, COLORS
\usepackage{graphics,graphicx,trimclip,adjustbox}     %figures and box management
\usepackage{color,xcolor}     %to use colors (used in the todonotes)
    
%%%%% ENUMERATE
\usepackage{enumitem}   %to use "enumerate"
    \setlist[enumerate,1]{label = $\bullet$}   %default label at bullet points
    
%%%%% TODO NOTES
\usepackage[textwidth=16.3mm,textsize=scriptsize]{todonotes}   %to put notes everywhere
\newcommand*{\newcommentator}[2][red!20]{%
    \expandafter\newcommand\csname #2\endcsname[1]{\todo[color=#1]{{\bf\makefirstuc{#2}:} ##1}}%
    \expandafter\newcommand\csname big#2\endcsname[1]{\todo[inline,color=#1]{{\bf\makefirstuc{#2}:} ##1}}}
\newcommentator{pablo}

%%%%% MY HYPERLINK AND REF/LABEL MACROS
\usepackage{mypreamble/linkref}

%%%%% MATH PACKAGES
\usepackage{amssymb,amsmath}     %math original package
    \numberwithin{equation}{section}
\usepackage{mathtools}     %improves AMSmath
\usepackage{stackrel}     %better stacking of math symbols
\usepackage{mleftright}     %defines \mleft and \mright
    \mleftright     %redefines \left as \mleft and \right as \mright
\usepackage{pgfcore}
\usepackage{tikz-cd}     %advanced and rigorous diagrams
    
%%%%% MY SPECIAL PACKAGES
\usepackage{mypreamble/mathenv}     %improves on AMSmath's equation
\usepackage{mypreamble/thmenv}     %replaces the AMSthm package

%%%%% MATH MACROS AND FONT PACKAGES
%\usepackage{fontspec}
\usepackage{mypreamble/mathcro}

%%%%% HYPERREF PACKAGE
\makeatletter
\pretocmd\refstepcounter{%
 \zref@setcurrent{counter}{#1}}{}{}
\makeatother
\usepackage[pdfencoding=auto]{hyperref}   %hyperlinks parts of the pdf
    \hypersetup{
        bookmarksopen=true,   %expends tree of bookmarks in Acrobat pdf reader
        bookmarksnumbered=true,  %show numbers of sections in bookmarks
        linktoc=all,   %both section and page numbers are links
        hidelinks   %links (color+border) are hidden
    }
\providecommand*{\definitionautorefname}{Definition}
\providecommand*{\notationautorefname}{Notation}
\providecommand*{\conventionautorefname}{Convention}
\providecommand*{\conjectureautorefname}{Conjecture}
\providecommand*{\assumptionautorefname}{Assumption}
\providecommand*{\remarkautorefname}{Remark}
\providecommand*{\observationautorefname}{Observation}
\providecommand*{\lemmaautorefname}{Lemma}
\providecommand*{\propositionautorefname}{Proposition}
\providecommand*{\corollaryautorefname}{Corollary}
\providecommand*{\exampleautorefname}{Example}
\providecommand*{\exerciseautorefname}{Exercise}
\providecommand*{\openquestionautorefname}{Open Question}

%%%%% OUTPUT LAYOUT PACKAGES
\usepackage{geometry}     %to set margins
\geometry{
    a4paper,
    total={170mm,257mm},
    left=20mm,
    top=20mm,
    marginparwidth=20mm,
}
%% Clearpage before sections but not before TOC
\let\oldsection\section
\renewcommand\section{\clearpage\oldsection}
\makeatletter
\renewcommand\tableofcontents{%
    \oldsection*{\contentsname\@mkboth{\MakeUppercase\contentsname}{\MakeUppercase\contentsname}}\@starttoc{toc}%
}
\makeatother    %ensures the TOC stays in front page

%%%%% TYPESET DEBUGGING
%\usepackage{showframe}     %to make margins appear
%\usepackage{lua-visual-debug}     %to make all boxes appear
%\overfullrule=1pt   %creates 1pt wide black bars next to any overfull hbox

%%%%% QUIVER SETUP
\usetikzlibrary{calc,decorations.pathmorphing}
% A TikZ style for curved arrows of a fixed height, due to AndréC.
\tikzset{curve/.style={settings={#1},to path={(\tikztostart)
    .. controls ($(\tikztostart)!\pv{pos}!(\tikztotarget)!\pv{height}!270:(\tikztotarget)$)
    and ($(\tikztostart)!1-\pv{pos}!(\tikztotarget)!\pv{height}!270:(\tikztotarget)$)
    .. (\tikztotarget)\tikztonodes}},
    settings/.code={\tikzset{quiver/.cd,#1}
        \def\pv##1{\pgfkeysvalueof{/tikz/quiver/##1}}},
    quiver/.cd,pos/.initial=0.35,height/.initial=0}

% TikZ arrowhead/tail styles.
\tikzset{tail reversed/.code={\pgfsetarrowsstart{tikzcd to}}}
\tikzset{2tail/.code={\pgfsetarrowsstart{Implies[reversed]}}}
\tikzset{2tail reversed/.code={\pgfsetarrowsstart{Implies}}}
% TikZ arrow styles.
\tikzset{no body/.style={/tikz/dash pattern=on 0 off 1mm}}

%%%%% TITLE PAGE
\title{Shifted Symplectic Integration}

\begin{document}

\maketitle

\section{First Definitions}

We work with $\infty$-stacks on the $\infty$-category of connective commutative derived rings in characteristic $0$: equivalently defined as the sind-cocompletion of polynomial algebra (free finitely generated $k$-algebras), simplicial commutative algebra, CDGA in positive (homological) degree, connective $\mathcal{E}^\infty$-algebras in $k\amod$,... The topology is étale and the stacks must have a suitable theory of (co)tangent modules (it has to be a dualizable object), for example we can use Artin stacks.

Ideally, I'd like to study stacks on $\dMan$ at some point: the point is that geometric quantization should be more likely to exist in the $\mathcal{C}^\infty$-setting rather than the algebraic one. It probably won't change any of what follows.

\begin{remark}[Working on Points]
    Recall for a stack $X \in \aff^{\op} \to \Spa$, we define
\[
    \qcoh(X) = \Lim_{X} (\aff^{\op} \xto[-\amod] \Cat)
\]
by point-wise right Kan extension. Automatically, this implies that $qcoh$ sends colimits of (pre)stacks to limits (of categories).

Let us consider $F\in\B \to \st$ a finite diagram of stacks. Then $\Lim F$ is a stack and we have a correspondence between $A$-points $x$ of $\Lim F$ and cones on $A$ of $F$. We can define functors
\[
    \mathbb{L}_x \in \begin{cases}[rCl]
        \mathcal{B} &\longrightarrow& A\amod\\
        b &\longmapsto& x_b^*(\mathbb{L}_{Fb})
    \end{cases}
\]
and its linear dual $\mathbb{T}_x$. We then have
\[
    x^*\mathbb{L}_{\Lim F} = \Colim_{\mathcal{B}} \mathbb{L}_x \\
    x^*\mathbb{T}_{\Lim F} = \Lim_{\mathcal{B}} \mathbb{T}_x
\]
Using this, building (shifted) maps between tangent and cotangent modules is equivalent to building such maps the respective (co)limits on the $A$-points, naturally in the points.
\end{remark}

\begin{remark}[How about \textbf{closed} forms?] It's not clear to me yet, how one could translate being a closed form at points like this? It's probably not possible: I will have to deal with this eventually.
\end{remark}

\section{Some Ideas}

What does the de Rham complex look like on a point?
Basically $d$ becomes $0$ but we still have symmetric powers of the cotangent at a point.
If we want to capture the behavior of $d$ we need to look at open covers, or something else.
The coefficient field is initial in CDGA and hence terminal in Aff. As a stack, it is also terminal. This gives a forgetful $\qcoh(X) \to k\amod$ and we can find $d$ in the image.

\section{Poincar\'e Duality}

Let $\B \to \Bsc$ be a tangential structure, \textit{i.e.} a right fibration. We have the following categories:
\[
    \Bsc(\mathcal{B})_+, \Bsc(\mathcal{B})^+ \longmono \ZBsc(\mathcal{B})
\]
where an object of $\Bsc(\mathcal{B})_+$ is a $\B$-basic adjointed a separate point $U\sqcup \{\infty\}$ and an object of $\Bsc(\mathcal{B})^+$ is a $\B$-basic adjointed a compactifying point $V^+$.
More explicitly, $\B \to \Bsc(\mathcal{B})_+$ canonically makes the hom-spaces pointed (by adding a discrete point), and $\Bsc(\mathcal{B})^+ \simeq \left(\Bsc(\mathcal{B})_+\right)^\op$ canonically.
In $\ZBsc(\mathcal{B})$, the hom-spaces between mixed types are either trivial or
\[
    \ZBsc(\mathcal{B})(U\sqcup \{\infty\}, V^+)
\]
which can be described by some simplicial set.

A functor $\B \to A\amod$ can be extended to $\Bsc(\mathcal{B})_+$ by using zero-maps and to $\Bsc(\mathcal{B})^+$ if it lands into dualizable objects.

\begin{definition}[Poincar\'e Context] A \textbf{Poincar\'e Context} is the data of an extension of a functor $\B \to A\amod$ to a Poincar\'e/Kozsul duality $\ZBsc(\mathcal{B}) \to A\amod$.
\end{definition}

\subsection{Examples}

\begin{enumerate}
    \item $\{\bullet\} \to \Bsc^{fr}$ picking $\mathbb{R}^n$ gives an object $A$ together with an isomorphism $A \simeq A^\vee[n]$
\end{enumerate}

\end{document}
